
% --------------------------------------------------
\subsubsection{Аналитик шугаман бус агшилтын үнэлгээ}
% --------------------------------------------------

Аналитик шугаман бус агшилтын үнэлгээ
(Analytical Nonlinear Shrinkage Estimation) нь
санамсаргүй матрицын онол (Random Matrix Theory)-д
тулгуурлан түүврийн ковариацын матрицын хувийн
утгуудыг засварлах арга юм. Шугаман агшилтын
аргуудаас ялгаатай нь энэхүү арга нь хувийн
утга бүрийн хэмжээ болон спектрийн тархалт дахь
байрлалыг харгалзан ялгаатай агшилт хийдэг.

\citet{ledoitwolf2020} нь спектрийн тархалтын
нягт болон түүний Hilbert хувиргалтыг kernel
аргаар үнэлэх замаар аналитик шугаман бус
агшилтын аргыг боловсруулсан. Уг арга нь
түүврийн ковариацын матрицын хувийн
векторуудыг хадгалж, хувийн утгуудыг
засварлах замаар ковариацын матрицын
үнэлгээний алдааг бууруулахад чиглэдэг.

Судалгаанд хамрагдсан хөрөнгийн тоог $N$,
ажиглалтын тоог $n$ гэж тэмдэглэж,
тэдгээрийн харьцааг
\[
c=\frac{N}{n}
\]
гэж тодорхойлно. Аргын үндсэн онолын хувьд
$N,n\to\infty$ үед
\[
\frac{N}{n}\longrightarrow c\in(0,1)
\]
байх асимптотик нөхцөлийг авч үздэг.

Түүврийн ковариацын матрицын спектрийн
задлалыг
\[
S=U\Lambda U^{\top}
\]
гэж тэмдэглэе. Энд
\[
U=(u_1,u_2,\ldots,u_N),
\qquad
U^{\top}U=I_N,
\]
нь хувийн векторуудаас бүрдэх ортогональ
матриц, харин
\[
\Lambda
=
\operatorname{diag}
(\lambda_1,\lambda_2,\ldots,\lambda_N)
\]
нь хувийн утгуудыг агуулсан диагональ
матриц юм. Хувийн утгуудыг
\[
\lambda_1\geq\lambda_2\geq\cdots
\geq\lambda_N>0
\]
гэж эрэмбэлнэ.

Аналитик шугаман бус агшилтын үнэлэгчийг
\[
\widehat{\Sigma}_{\mathrm{NLS}}
=
U\widehat{D}U^{\top}
=
\sum_{i=1}^{N}
\widehat{d}_i u_i u_i^{\top}
\]
гэж тодорхойлно. Энд
\[
\widehat{D}
=
\operatorname{diag}
(\widehat{d}_1,\ldots,\widehat{d}_N)
\]
нь агшсан хувийн утгуудаас бүрдэх
диагональ матриц юм.

\paragraph{Агшсан хувийн утгуудын үнэлгээ.}

\citet{ledoitwolf2020}-ийн (4.3) томьёоны
дагуу агшсан хувийн утгуудыг
\[
\widehat{d}_i
=
\frac{\lambda_i}
{
	\left(
	1-c-c\pi\lambda_i
	\widehat{Hf}(\lambda_i)
	\right)^2
	+
	\left(
	c\pi\lambda_i
	\widehat{f}(\lambda_i)
	\right)^2
}
\]
гэж тооцно.

Энд $\widehat{f}(\lambda_i)$ нь
спектрийн тархалтын нягтын үнэлгээ,
$\widehat{Hf}(\lambda_i)$ нь уг нягтын
Hilbert хувиргалтын үнэлгээ юм.
Эдгээрийг түүврийн ковариацын матрицын
хувийн утгуудад тулгуурлан kernel
аргаар тооцдог.

Иймээс агшсан хувийн утга нь зөвхөн
тухайн $\lambda_i$-ийн хэмжээнээс бус,
спектрийн тархалтын нягт болон түүний
Hilbert хувиргалтын утгаас хамаарна.

\paragraph{Тооцоолох алгоритм.}

\citet{ledoitwolf2020}-ийн Section 4.7-д
үндэслэн аналитик шугаман бус агшилтын
ковариацын матрицыг дараах үе шаттайгаар
тооцно.

\begin{enumerate}
	

	\item Түүврийн ковариацын матрицын
	спектрийн задлалыг байгуулна.
	\[
	S=U\Lambda U^{\top}.
	\]
	
	\item Хөрөнгийн тоо болон ажиглалтын
	тооны харьцааг тооцно.
	\[
	c=\frac{N}{n}.
	\]
	
	\item Спектрийн нягтыг үнэлэхэд
	ашиглах bandwidth параметрийг
	\[
	h=n^{-1/3}
	\]
	гэж сонгоно. Хувийн утга бүрд
	харгалзах локал bandwidth нь
	\[
	h_j=\lambda_jh,
	\qquad j=1,\ldots,N
	\]
	байна.
	
	\item Спектрийн нягтыг Epanechnikov
	kernel ашиглан үнэлнэ.
	\[
	\widehat{f}(\lambda_i)
	=
	\frac{1}{N}
	\sum_{j=1}^{N}
	\frac{1}{h_j}
	k\left(
	\frac{\lambda_i-\lambda_j}{h_j}
	\right),
	\]
	энд
	\[
	k(x)
	=
	\begin{cases}
		\dfrac{3}{4\sqrt{5}}
		\left(1-\dfrac{x^2}{5}\right),
		& |x|\leq\sqrt{5},\\[0.5em]
		0,
		& |x|>\sqrt{5}.
	\end{cases}
	\]
	
	\item Спектрийн нягтын Hilbert
	хувиргалтыг үнэлнэ.
	\[
	\widehat{Hf}(\lambda_i)
	=
	\frac{1}{N}
	\sum_{j=1}^{N}
	\frac{1}{h_j}
	\ell\left(
	\frac{\lambda_i-\lambda_j}{h_j}
	\right),
	\]
	энд $\ell(x)=Hk(x)$ бөгөөд
	\[
	\ell(x)
	=
	\begin{cases}
		-\dfrac{3x}{10\pi}
		+
		\dfrac{3}{4\sqrt{5}\pi}
		\left(1-\dfrac{x^2}{5}\right)
		\ln\left|
		\dfrac{\sqrt{5}-x}{\sqrt{5}+x}
		\right|,
		& |x|\neq\sqrt{5},\\[1em]
		-\dfrac{3x}{10\pi},
		& |x|=\sqrt{5}.
	\end{cases}
	\]
	
	\item Агшсан хувийн утгуудыг
	тооцно.
	\[
	\widehat{d}_i
	=
	\frac{\lambda_i}
	{
		\left(
		1-c-c\pi\lambda_i
		\widehat{Hf}(\lambda_i)
		\right)^2
		+
		\left(
		c\pi\lambda_i
		\widehat{f}(\lambda_i)
		\right)^2
	}.
	\]
	
	\item Агшилтын ковариацын матрицыг
	байгуулна.
	\[
	\widehat{\Sigma}_{\mathrm{NLS}}
	=
	U
	\operatorname{diag}
	(\widehat{d}_1,\ldots,\widehat{d}_N)
	U^{\top}.
	\]
	
\end{enumerate}
